% !TeX program = xelatex
\documentclass[11pt]{article}

\usepackage[a4paper, margin=1in]{geometry}
\usepackage{amsmath, amssymb, amsthm}
\usepackage{stmaryrd}
\usepackage{listings}
\usepackage{xcolor}
\usepackage{hyperref}
\usepackage{microtype}
\usepackage{fontspec}
\usepackage{booktabs}
\usepackage{longtable}
% Use DejaVu Sans Mono if present (as the other Notes files do); fall back otherwise.
\IfFontExistsTF{DejaVu Sans Mono}{\setmonofont{DejaVu Sans Mono}[Scale=0.85]}{}

\sloppy
\setlength{\emergencystretch}{3em}

% Robust so it survives moving arguments (section titles -> table of contents).
\DeclareRobustCommand{\btt}[1]{\texttt{\detokenize{#1}}}

% ---- Lean code listing style -------------------------------------------
\lstdefinelanguage{lean}{
  morekeywords={theorem, lemma, def, by, fun, let, in, do, match, with,
    if, then, else, return, import, open, namespace, end, structure,
    inductive, partial, private, noncomputable, syntax, elab, elab_rules,
    macro, macro_rules, tactic, conv, command, example, have, show, intro,
    apply, simp, rw, ring, all_goals, repeat, decide, native_decide,
    fin_cases, norm_num, exact, set_option, unfold, induction},
  sensitive=true,
  morecomment=[l]{--},
  morecomment=[s]{/-}{-/},
  morestring=[b]"
}
\lstset{
  inputencoding=utf8,
  language=lean,
  basicstyle=\ttfamily\footnotesize,
  keywordstyle=\color{blue!70!black}\bfseries,
  commentstyle=\color{gray}\itshape,
  stringstyle=\color{red!60!black},
  showstringspaces=false,
  breaklines=true,
  columns=fullflexible,
  frame=single,
  framesep=4pt,
  rulecolor=\color{black!30},
  backgroundcolor=\color{gray!6},
  xleftmargin=6pt,
  xrightmargin=6pt,
  aboveskip=0.8\baselineskip,
  belowskip=0.8\baselineskip
}
% -----------------------------------------------------------------------

\title{Generating and Optimizing Flags in Lean:\\
An Implementation and Optimization History of\\
\btt{generate_empty_typed_flags} and \btt{generate_flags}}
\author{}
\date{}

\begin{document}
\maketitle

\begin{abstract}
This note documents the design, implementation, and optimization of the two
elaboration-time flag-generation commands of the \btt{lean-flag-algebras}
development, \btt{generate_empty_typed_flags} and \btt{generate_flags}. It traces
the code from commit \btt{afe184a} (``Phase~1''), where the empty-typed flags were
already generated in Lean but the typed (labeled) flags were still loaded from
external JSON, through the cut-over to fully self-contained Lean generation and a
long sequence of performance and proof-engineering optimizations, ending with the
augmentation-based labeled enumeration and the degree-key dedup prefilter. For each
change we explain how it worked, why it was made, what improvement it produced (with
measured numbers wherever they were recorded), and what remained slow afterwards. We
also document the cases where an anticipated improvement failed to materialize, the
dead ends and confounded diagnoses, the two fundamental walls that make $n=6$
expensive and $n=7$ impractical under the current architecture, and the directions
that remain promising. The note is intended to be self-contained for a reader who has
never seen this code.
\end{abstract}

\tableofcontents

% =======================================================================
\section{Scope and orientation}
% =======================================================================

The \btt{lean-flag-algebras} project formalizes Razborov-style flag-algebra
arguments (here, the Murphy--Nir 2021 results). A flag-algebra proof works over a
\emph{basis of flags}: small graphs, possibly carrying a labeled ``type''~$\sigma$.
To run such a proof inside Lean one must \emph{enumerate} all flags of a given size
(one canonical representative per isomorphism class) and \emph{prove} two facts about
the enumeration:
\begin{itemize}
  \item \textbf{Completeness}: the enumeration hits every isomorphism class, i.e.\
        its image in the quotient is \btt{Finset.univ}.
  \item \textbf{No duplication}: \btt{Nodup} of the enumerated list.
\end{itemize}
Originally these flags were produced by external Python scripts that emitted JSON,
which a Lean macro parsed at elaboration time. The work documented here replaced that
JSON pipeline with self-contained Lean enumeration: two \emph{commands} that, at
elaboration time, run a pure Lean computation to enumerate the flags and synthesize
the named constants and completeness/no-duplication/downward-coefficient theorems that
the rest of the development consumes.

The two commands are:
\begin{itemize}
  \item \btt{generate_empty_typed_flags n} --- the unlabeled (``empty-typed'',
        type $\varnothing_t$) flags on $n$ vertices: one canonical graph per
        isomorphism class, together with the abstract \btt{Flag} bridges.
  \item \btt{generate_flags k m n} --- the $\sigma$-typed (labeled) flags on $n$
        vertices, where $\sigma$ is the $k$-vertex graph of canonical index $m$.
        Each labeled flag is built on top of an underlying empty-typed flag and
        additionally records its \emph{type embedding} and a downward-normalizing
        coefficient.
\end{itemize}

This document reviews every relevant commit from \btt{afe184a} (2026-06-04) to
\btt{54b06f7}, the last commit at the time of writing. Throughout, $g = g(n,k,m)$
denotes the number of flags produced by a line (the count of isomorphism classes),
which is the dominant scaling parameter. Representative counts: the empty-typed
classes are $1,1,2,4,11,34,156,1044$ for $n=0,\dots,7$; the heaviest typed lines used
in practice are $\btt{generate_flags 3 0 5}$ ($g=72$) and $\btt{generate_flags 3 0 6}$
($g=816$).

% =======================================================================
\section{Background: the data model and the generation paradigm}
% =======================================================================

\subsection{Computable graphs, labeled graphs, and flags}

The development carries two parallel representations of the same mathematical
objects. The abstract layer (\btt{LabeledGraph}, \btt{Flag}) is the one the
flag-algebra theory is stated over; the \emph{computable} layer (\btt{Sym2Graph},
\btt{Sym2LabeledGraph}) is finite-data and decidable, suitable for
\btt{decide}/\btt{native_decide}.

\begin{itemize}
  \item \btt{Sym2Graph n} is a graph on \btt{Fin n}: a \btt{Finset (Sym2 (Fin n))}
        of edges together with a proof that no edge is a self-loop
        (\btt{edges_valid}). \btt{Sym2EmptyTypedFlag n} is the quotient of
        \btt{Sym2Graph n} by isomorphism; \btt{Sym2Graph.toFlag} maps it into the
        abstract \btt{Flag} $\varnothing_t\ (\btt{Fin n})$.
  \item \btt{Sym2LabeledGraph} $\sigma$ \btt{n} extends a \btt{Sym2Graph} with a
        \emph{type embedding} \btt{type_embed}, a graph embedding of (the decoded)
        $\sigma$ into the underlying graph. \btt{Sym2Flag} $\sigma$ \btt{n} is its
        quotient by labeled isomorphism, and \btt{Sym2Flag.toFlag} maps it to the
        abstract \btt{Flag}.
  \item The computable isomorphism relation is written $\sim_{\!sf}$
        (\btt{Sym2GraphEqv} / \btt{sym2LabeledGraphEqv}); two computable
        (labeled) graphs are equivalent iff their decoded abstract forms are
        flag-isomorphic. The fast decidable check is \btt{isEmptyIsoFast_bool}
        (unlabeled) / \btt{isIsoFast_bool} (labeled).
\end{itemize}

A \emph{flag isomorphism} of labeled graphs carries a \btt{graph_iso} (an isomorphism
of the underlying simple graphs) together with a \btt{type_preserve} field stating
that the isomorphism commutes with the two type embeddings. This structure is what
makes the labeled iso check both subtler and, as it happens, \emph{cheaper} than the
unlabeled one (Section~\ref{sec:degkey}).

\subsection{The two stages of a generation command}

Each command has two stages, and almost every optimization in this history targets one
of them:

\begin{enumerate}
  \item \textbf{Elaboration-time data computation.} The macro evaluates a pure Lean
        function (\btt{genCanonicalEdgeLists n} or \btt{genFlagData k m n}) to obtain
        the flag data --- canonical edge lists, type indices, and downward
        coefficients. This evaluation runs in the Lean elaborator/interpreter (or via
        compiled \btt{Lean.Meta.evalExpr}).
  \item \textbf{Declaration emission and proof.} For each flag the macro emits named
        \btt{def}s (\btt{Sym2Graph_n_0_0_i}, \btt{Sym2Flag_n_k_m_i},
        \btt{Flag_n_k_m_i}, \btt{FlagAlgebra_n_k_m_i}, and for labeled flags an
        \btt{unlabel_n_k_m_i} bridge to the underlying flag and a
        \btt{downward_n_k_m_i} coefficient theorem), and finally the set-level
        machinery: a finset of the named flags, a \btt{Nodup} proof, and a
        \btt{= Finset.univ} completeness proof.
\end{enumerate}

\subsection{The central proof problem: completeness without the ``monster''}
\label{sec:monster}

The naive way to prove \btt{sym2FlagSet_n_k_m = Finset.univ} is a single
\btt{native_decide}. But \btt{Finset.univ : Finset (Sym2Flag} $\sigma$ \btt{n)}
forces the \btt{Fintype (Sym2LabeledGraph} $\sigma$ \btt{n)} instance, which
enumerates \emph{all} $2^{\binom n2}$ edge subsets crossed with \emph{all} type
embeddings, then quotients --- an astronomically large object we call the
\emph{monster}. Materializing it is the original bottleneck.

The architecture's answer, common to both commands, is a \emph{symbolic} completeness
theorem proved once, mathematically, in \btt{FlagEnumeration.lean}:
\begin{itemize}
  \item \btt{genEmptyTypedFlagSet_eq_univ n} and
        \btt{genFlagSet_eq_univ}~$\sigma$~\btt{n}: the \emph{generated} finset equals
        \btt{univ}, proved by \btt{Finset.eq_univ_of_forall} from a completeness lemma
        about the generator --- never by reduction.
\end{itemize}
The macro then only has to \emph{bridge} the concrete named-flag list to the
generator's enumeration. That bridge is the part discharged by \btt{native_decide},
and reducing its cost from $O(g^2)$ to $O(g)$ is the recurring theme of this history.

% =======================================================================
\section{The original implementation and the \btt{afe184a} baseline}
% =======================================================================

\subsection{The JSON pipeline (pre-Phase~1)}

In the original implementation, two external Python scripts
(\btt{generate_graphs.py}, \btt{generate_flags.py}) enumerated the
non-isomorphic graphs and the typed flags and wrote them to JSON files
(\btt{graphs_n.json}, \btt{flags_n_k_m.json}). A Lean macro module (then named
\btt{FlagLoader.lean}) provided \btt{load_empty_typed_flags "graphs_n.json"} and
\btt{load_flags "flags_n_k_m.json"}, which parsed the JSON with
\btt{Lean.Data.Json} and emitted the corresponding named declarations. The
provenance of the enumeration was thus \emph{outside} the Lean development: a reader
had to trust the Python scripts and the JSON files, and the build depended on those
files being present and consistent.

\subsection{Phase~1 (commit \btt{afe184a}): empty-typed flags in Lean}

At \btt{afe184a} (``Phase~1: generate empty-typed flags in Lean, replacing the JSON
pipeline'') the \emph{empty-typed} half had already been internalized.
\btt{FlagEnumeration.lean} contained a pure Lean enumeration of one canonical graph per
isomorphism class, and \btt{generate_empty_typed_flags n} emitted those graphs and
their flags. The key pieces, which all later work builds on, were:
\begin{itemize}
  \item \textbf{Canonical-augmentation generator.} \btt{augReps n} builds the
        $n$-vertex representatives by induction on $n$: take the $n$-vertex reps, add
        one new vertex adjacent to each subset of the existing vertices
        (\btt{augmentAll}), and deduplicate up to $\sim_{\!sf}$ by a fold
        (\btt{dedupStep}). This is the standard ``orderly/canonical augmentation''
        idea; crucially it never enumerates all $2^{\binom n2}$ subsets --- it grows
        the representatives directly.
  \item \textbf{Canonicalization and ordering.} \btt{canonicalEdgeList} computes a
        graph's lexicographically minimal edge list over all $n!$ relabelings (the
        canonical form); \btt{graphKey} $=$ (edge count, canonical edge list)
        provides a total order; the generated list \btt{genSym2Graphs n} is sorted
        into this canonical/``JSON'' order so the named constants come out in a
        stable, reproducible order.
  \item \textbf{Completeness and no-duplication}, proved symbolically:
        \btt{augReps_complete} (every graph is $\sim_{\!sf}$ to a representative),
        \btt{genEmptyTypedFlags_nodup}, and \btt{genEmptyTypedFlagSet_eq_univ}.
\end{itemize}
At \btt{afe184a}, \btt{FlagDef.lean} generated empty-typed flags for $n=0,\dots,5$ in
Lean, but still loaded the typed flags from JSON via \btt{load_flags}, and the JSON
\btt{load_*} commands still existed alongside the generators. This is our baseline.

% =======================================================================
\section{Phase~2: typed-flag generation and symbolic completeness (\btt{212968a})}
% =======================================================================

\subsection{What changed}

Commit \btt{212968a} (``Phase~2'') internalized the typed flags. It added
\btt{genFlagData k m n} to \btt{FlagEnumeration.lean}: for each underlying graph (in the
canonical order of \btt{genSym2Graphs n}) it enumerates the valid type embeddings of
$\sigma$, groups them into orbits under the underlying graph's automorphism group,
keeps the orbit-lexicographically-minimal tuple as the representative, and records the
downward coefficient $\lvert\text{orbit}\rvert / n^{\underline{k}}$ in lowest terms.
The new command \btt{generate_flags k m n} consumes this data, emitting one
\btt{Sym2LabeledGraph}/\btt{Sym2Flag}/\btt{Flag}/\btt{FlagAlgebra} per flag, plus the
per-flag \btt{unlabel} and \btt{downward} bridges, plus the set-level machinery.

\subsection{Why, and the improvement}

The point was to remove the JSON dependency \emph{and} to avoid the monster. Instead
of a \btt{native_decide} over \btt{Finset.univ : Finset (Sym2Flag} $\sigma$ \btt{n)},
\btt{sym2FlagSet_n_k_m_eq_univ} now applies the symbolic \btt{genFlagSet_eq_univ}
(Section~\ref{sec:monster}), bridged to the named list by one \btt{native_decide} over
the \emph{tight} \btt{genFlagSet} enumeration. The recorded effect was a drop of the
\btt{FlagDef} build from \textbf{568\,s to 274\,s} ($\sim 2\times$), with all flag
counts unchanged and the full project compiling cleanly. The bridge at this point was
still an $O(g^2)$ comparison (a finset/quotient equality, each of the $g$ elements
membership-checked, each membership an iso check), which the later optimizations
target.

Commit \btt{8a3ceab} then removed the now-dead JSON loaders
(\btt{load_empty_typed_flags}, \btt{load_flags}, all the parsing helpers and JSON
structs, and the \btt{Lean.Data.Json} import) --- roughly 640 lines of dead code ---
completing the cut-over.

% =======================================================================
\section{Empty-typed optimizations (\btt{fbaeed0}, \btt{fa8939e}, \btt{de7dcbc})}
% =======================================================================

These three commits sharpened the empty-typed path and pushed it to $n=6$; several of
their ideas were later ported to the labeled path during the session work.

\subsection{Degree-key prefilter, decorate-sort-undecorate, and the $O(g)$ bridge (\btt{fbaeed0})}
\label{sec:fbaeed0}

This commit made three orthogonal changes:

\paragraph{Degree-sequence dedup prefilter.} The augmentation dedup compared each new
candidate against every survivor with the $O(n!)$-flavored \btt{isEmptyIsoFast_bool}.
\btt{fbaeed0} introduced a cheap, $O(n^2)$ isomorphism invariant
\btt{degKey} $=$ (edge count, sorted degree sequence) and a keyed dedup step
\btt{dedupStepDeg} that runs the expensive iso test \emph{only} when the cheap keys
collide:
\begin{lstlisting}
def dedupStepDeg (acc) (p) :=
  if acc.any (fun q => q.2 == p.2 && isEmptyIsoFast_bool q.1 p.1) = true
  then acc else acc ++ [p]
\end{lstlisting}
Correctness is non-trivial: the keyed fold must make the \emph{same} keep/drop
decisions as the naive fold. This needs (i) \btt{degKey_iso_invariant}: $\sim_{\!sf}$
graphs have equal keys (so the prefilter never wrongly skips a real isomorph), proved
via \btt{degree_eq_of_eqv}/\btt{degSeq_eq_of_eqv}; and (ii) a fold simulation
\btt{foldl_dedupStepDeg_sim} carrying two invariants (the first components match, every
tag is its graph's \btt{degKey}). Finally \btt{augRepsDeg_fst_eq} shows the fast keyed
generator's first components equal the specification \btt{augReps}, so all
completeness/\btt{Nodup} reasoning continues to run on \btt{augReps}. The whole change
is \btt{sorry}-free.

\paragraph{Decorate-sort-undecorate.} \btt{canonicalEdgeList} is $O(n!)$. Computing it
afresh inside the sort comparator would cost $O(g^2)$ key computations. The keyed
generator instead \emph{decorates} each graph with its key once, sorts by the
precomputed key, and carries the key through --- $O(g)$ key computations plus a cheap
$O(g^2 \log g)$ comparison sort on small tuples.

\paragraph{$O(g)$ completeness bridge.} The empty-typed completeness bridge was
rewired from an $O(g^2)$ \btt{toFinset} comparison to a \emph{positional list
equality}
\[
  \btt{[Sym2Flag_n_0_0_i] = genEmptyTypedFlags n},
\]
which costs $O(g)$ iso checks --- one per position --- because the named list and the
generator are in the \emph{same} canonical order (this is exactly what the sort layer
buys). This is the template the labeled path eventually copied (commit \btt{bb6f60f}).
The commit also extended \btt{FlagDef} to $n=6$ (156 classes) and documented why
$n=7$ (1044) is omitted under interpreter-based generation.

\subsection{Structural \btt{edges_valid} (\btt{fa8939e})}
\label{sec:edges_valid}

Every generated graph carried an \btt{edges_valid} proof (no edge is a self-loop). It
was discharged by \btt{by decide} over \btt{l.toFinset}, which forces the kernel to
re-run the full \btt{Decidable} reduction to check \btt{of_decide_eq_true (Eq.refl
true)} --- paid once in the elaborator and again in kernel type-checking, for every
graph. The fix added \btt{mkEdgeFinset_diag_free}, reducing finset-level off-diagonal
validity to a list check (membership in \btt{l.toFinset} is membership in \btt{l}),
discharged structurally by \btt{fin_cases <;> simp [Sym2.isDiag_iff_proj_eq]}. The
proof term becomes compact and cheap for the kernel. On the $n=6$ worst case (156
graphs) this cut \btt{decide} tactic execution $\sim 9\times$ and roughly halved kernel
type-checking for these definitions.

\subsection{Reusing the precomputed canonical edge list (\btt{de7dcbc})}
\label{sec:reuse}

\btt{genSym2Graphs} decorated each graph with its \btt{graphKey} $=$ (edge count,
canonical edge list), sorted, then \emph{dropped} the keys via \btt{.map Prod.fst}.
Four downstream sites (the two generation commands and both lines of
\btt{genFlagData}) then \emph{recomputed} the $O(n!)$ \btt{canonicalEdgeList} from
scratch. \btt{de7dcbc} threads the keys through instead: \btt{genSym2GraphsKeyed}
keeps each graph paired with its key, \btt{genSym2Graphs} becomes a projection, and
\btt{genCanonicalEdgeLists} reads the canonical lists straight off the keys (value-
identical because the insertion sort permutes whole pairs, so $p.2.2 =
\btt{canonicalEdgeList}\ p.1$ is invariant). This is a pure value-preserving refactor:
\btt{genSym2Graphs_perm} only needed \btt{genSym2GraphsKeyed} added to its
\btt{unfold}. On $n=6$ the empty-typed generation interpretation dropped
\textbf{13.6\,s $\to$ 9.73\,s ($\sim 28\%$)}.

% =======================================================================
\section{Typed-flag elaboration cost: whole-goal \btt{decide} (\btt{81cf849})}
% =======================================================================

The generated \btt{type_embed} field carries two proofs per flag: that the embedding
map is injective, and an adjacency-preservation statement \btt{hmap}. These were
discharged by \btt{fin_cases ... <;> decide}/\btt{simp}, which re-synthesizes a
\btt{Decidable}/\btt{Fintype} instance \emph{per case per flag} --- a typeclass-
inference blow-up. Discharging each whole goal with a single \btt{decide} synthesizes
the instance once. The heaviest typed line ($n=5$, $k=3$, 72 flags) dropped
\textbf{$\sim 63$\,s $\to \sim 22$\,s}, and the full \btt{FlagDef} build
\textbf{$\sim 301$\,s $\to \sim 100$\,s ($\sim 3\times$)}, with all counts and
downstream theorems unchanged. This was the single largest end-to-end win in the
typed path and the reason the per-flag \btt{decide}s are no longer the dominant cost.

% =======================================================================
\section{Refactors and naming (\btt{4b6632d}, \btt{bc53254}, \btt{36a6741})}
% =======================================================================

Three maintainability commits, with no intended performance effect:
\btt{4b6632d} replaced the per-declaration \btt{getEnv}/\btt{env.contains} idempotence
guard repeated at every emission with a single \btt{elabUnlessDefined} helper, and
collapsed the near-verbatim finset-to-\btt{Flag} bridge trio (\btt{flagSet},
\btt{_val_eq}, \btt{_eq_univ}) duplicated across both macros into one parameterized
\btt{emitFlagSetMachinery} emitter. \btt{bc53254} standardized a finset name.
\btt{36a6741} renamed \btt{FlagLoader.lean} to \btt{FlagGenerator.lean}, since the
module no longer loads anything --- its macros synthesize flags from a self-contained
Lean enumeration. These refactors matter for the rest of the story only because they
created the shared helpers (\btt{elabUnlessDefined}, \btt{emitFlagSetMachinery}) that
the session work plugged into.

% =======================================================================
\section{Session work I: the $O(g)$ labeled completeness bridge (\btt{bb6f60f})}
% =======================================================================

\subsection{The starting asymmetry}

After the pre-session work, the empty-typed path had an $O(g)$ positional bridge
(Section~\ref{sec:fbaeed0}) but the labeled path did \emph{not}: \btt{generate_flags}
still discharged completeness by an $O(g^2)$ \btt{Finset}-equality
\btt{sym2FlagSet = genFlagSet}, \emph{and} a \emph{separate} $O(g^2)$ \btt{Nodup}
\btt{native_decide}. The session began by closing this gap.

\subsection{Prototype, then measure, then prove}

The labeled enumeration was not sorted into canonical order, so a positional bridge
needed a sort layer. The approach mirrored \btt{genSym2GraphsKeyed}: define
\btt{genFlagsOrdered} $=$ the deduped reps re-sorted into \btt{genFlagData}'s order via
a sort key \btt{labeledFlagKey} (the underlying graph's \btt{graphKey} together with
the orbit-minimal type tuple in canonical coordinates). Crucially, the key is built
once per representative and carried through the sort (the same decorate-sort-undecorate
discipline; a first attempt that recomputed the $n!$-scanning key inside the comparator
ballooned the cost and was discarded).

A deliberate decision was to \emph{prototype and measure before proving}: the key and
\btt{genFlagsOrdered} were added as definitions, the bridge
\btt{[flagTerms] = genFlagsOrdered} was checked to evaluate \btt{true} by
\btt{native_decide} (confirming the key reproduces the JSON order), and the per-bridge
costs were measured in isolation before any perm/\btt{Nodup} proof was written. The
measurement at $n=5$ ($g=72$), as three separately profiled \btt{native_decide}s:

\begin{center}
\begin{tabular}{lll}
\toprule
bridge & proves & kernel cost \\
\midrule
\btt{measO2} (old eq\_univ) & \btt{sym2FlagSet = genFlagSet} ($O(g^2)$) & 4.5\,s \\
\btt{measNodup} (old nodup) & \btt{[flagTerms].Nodup} ($O(g^2)$)       & 13.1\,s \\
\btt{measO1} (new bridge)   & \btt{[flagTerms] = genFlagsOrdered} ($O(g)$) & 7.7\,s \\
\bottomrule
\end{tabular}
\end{center}

\subsection{A false alarm, then the real comparison}

Read naively, the new bridge ($7.7$\,s) looked \emph{slower} than the old eq\_univ
bridge ($4.5$\,s) --- a near-miss that would have killed the change. The resolution was
that this is the wrong comparison. The \emph{old} design pays \emph{two} $O(g^2)$
\btt{native_decide}s per line, \btt{measO2} \emph{and} \btt{measNodup}, because
\btt{Nodup} of the quotient list is itself $O(g^2)$ and, worse, every one of its
$\sim g^2/2$ pairwise checks is a \emph{failed} iso search (a full permutation scan to
conclude non-equality) --- which is why \btt{measNodup} is the single most expensive
piece at $13.1$\,s. The \emph{new} design pays only \btt{measO1}; both \btt{eq_univ}
and \btt{Nodup} are then derived from that one bridge by ordinary proof, with no further
\btt{native_decide}. So the honest per-line comparison is
\[
  \underbrace{4.5 + 13.1 = 17.6\,\text{s}}_{\text{old}} \;\longrightarrow\;
  \underbrace{7.7\,\text{s}}_{\text{new}} \quad (\sim 2.3\times),
\]
and the gap widens at $n=6$ because the $O(g^2)$ pieces grow faster than the $O(g)$
one. A decomposition baseline showed a shared $\sim 3.2$\,s dedup-fold cost inside both
\btt{measO2} and \btt{measO1}, so the genuine bridge-comparison delta is even more
favorable.

\subsection{The proofs}

The supporting math, all added to \btt{FlagEnumeration.lean}, was small and low-risk because
it mirrors the empty-typed template:
\begin{itemize}
  \item \btt{genFlagsOrdered_perm}: \btt{genFlagsOrdered} $\sim$ \btt{genFlags} (a
        permutation), a $\sim$12-line near-copy of \btt{genSym2Graphs_perm} using
        \btt{List.perm_insertionSort}.
  \item \btt{genFlags_nodup} and the helper \btt{foldl_dedupStepL_flags_nodup}: a
        structural induction on the dedup fold, \emph{no} \btt{native_decide}.
  \item \btt{genFlagsOrdered_nodup}, \btt{genFlagsOrdered_toFinset}: transfer
        \btt{Nodup} and \btt{= genFlagSet} across the permutation.
\end{itemize}
The macro was then rewired to emit one bridge theorem \btt{Sym2FlagList_n_k_m_eq}
($O(g)$ \btt{native_decide}) and to derive \btt{eq_univ} (through the order-independent
\btt{genFlagSet_eq_univ}) and the \btt{Nodup} obligation (through \btt{genFlags_nodup})
by proof. Notably, the sort key itself needs \emph{no} correctness proof: it lives in
the \btt{native_decide}-trusted zone, exactly like the JSON order it replaces --- if
it were wrong, the bridge simply would not compile. Verified: all 16 \btt{generate_flags}
lines through $n=5$ build, counts unchanged, full project green.

% =======================================================================
\section{Session work II: augmentation-based labeled enumeration (\btt{a993a33})}
% =======================================================================

\subsection{The exponential the labeled path was still paying}
\label{sec:bruteforce}

The $O(g)$ bridge above made the \emph{comparison} cheap, but it did not change
\emph{what the bridge evaluates}. The bridge runs \btt{native_decide} on
\btt{genFlagsOrdered}, which sits on \btt{genLabeledGraphsDedup}, which deduplicated
\begin{lstlisting}
allRawSym2LabeledGraphs sigma n = (allRawSym2Graphs n).flatMap (labeledOfGraph sigma)
\end{lstlisting}
--- i.e.\ \emph{every} one of the $2^{\binom n2}$ edge subsets of $K_n$ crossed with
each type embedding. At $n=6$ this is \textbf{491{,}520} raw labeled graphs
($2^{15}=32768$ underlying subsets $\times$ embeddings), and the dedup fold runs on the
order of \textbf{$10^8$} isomorphism checks. This brute force was \emph{exactly} the
exponential sweep the empty-typed path had already escaped via canonical augmentation;
the labeled path had simply never been given the same treatment. This was the dominant
$n=6$ cost and the subject of an explicit user decision (``go with option~2'').

\subsection{The fix and its one new proof}

The change replaces the brute-force factor \btt{allRawSym2Graphs n} (32768 graphs at
$n=6$) with the augmentation-generated canonical graphs \btt{genSym2Graphs n} (156 at
$n=6$):
\begin{lstlisting}
def allAugSym2LabeledGraphs sigma n := (genSym2Graphs n).flatMap (labeledOfGraph sigma)
\end{lstlisting}
The subtlety is completeness. The old chain used \emph{literal} membership
(\btt{mem_allRawSym2LabeledGraphs}: every labeled graph is \emph{in} the brute-force
list, since all subsets are present). The new list does not literally contain every
labeled graph, so literal membership is replaced by \emph{up-to-iso} completeness:
\begin{itemize}
  \item \btt{mem_labeledOfGraph_eqv_of_underlying} (the new content): if a labeled
        graph $H$'s underlying graph is $\sim_{\!sf}$ to a canonical $R$, then $H$ is
        $\sim_{\!sf}$ to a labeled graph over $R$ --- obtained by \emph{transporting}
        $H$'s type embedding along the underlying-graph isomorphism (post-composing
        with the iso $\psi$). The flag-equivalence witness is
        \btt{{ graph_iso := psi, type_preserve := rfl }}: because the new embedding is
        \emph{defined} as $\psi \circ H.\btt{type_embed}$, the \btt{type_preserve}
        obligation is literally \btt{rfl}, exactly as the design predicted.
  \item \btt{allAugSym2LabeledGraphs_complete}: combine
        \btt{genSym2Graphs_complete} (underlying completeness, already proved on the
        empty-typed side) with the transport.
\end{itemize}
Everything downstream is unchanged: \btt{genFlagsOrdered} sorts by an iso-invariant
key, so its \emph{value} (and hence the bridge \btt{native_decide}, and the flag
counts) is provably identical regardless of which representatives the dedup picks. The
$n=5$ build confirmed identical counts and that every bridge still evaluates
\btt{true}.

\subsection{Measured improvement}

The fast enumeration is \textbf{2{,}808} raw labeled graphs at $n=6$ versus
\textbf{491{,}520} --- a \textbf{$175\times$} reduction --- and the dedup fold drops from
$\sim 10^8$ to $\sim 2.3\times 10^6$ isomorphism checks. The whole computation (with the
full dedup) runs in well under a minute compiled. This removed the \emph{enumeration}
wall completely; what remained at $n=6$ is a different, smaller set of costs
(Section~\ref{sec:n6frontier}).

% =======================================================================
\section{The $n=6$ frontier and the missing-prerequisite bug (\btt{b15c62b})}
\label{sec:n6frontier}
% =======================================================================

With the enumeration cheap, an attempt to actually build \btt{generate_flags 3 0 6}
ran into a cascade of \emph{distinct} walls, plus one outright bug. This section is
also a record of how the diagnosis went wrong before it went right, because the user
asked for the dead ends.

\subsection{Three stacked resource walls}

\begin{enumerate}
  \item \textbf{\btt{maxRecDepth} (512).} \btt{generate_flags 3 0 6} emits list
        literals of length $g=816$ (the bridge \btt{[flagTerms]}, the \btt{setName}
        \btt{toFinset}, the downward batch \btt{[dnfTerms]/[coeffTerms]}). Each is an
        816-deep \btt{List.cons} term; elaborating and kernel-reducing it exceeds the
        default \btt{maxRecDepth} of 512. (This is the same wall the $n=7$ empty-typed
        case is documented to hit at 1044 elements.) Confirmed by \btt{\#eval}:
        $g(3,0,6)=816 > 512$. Raising \btt{maxRecDepth} to $\sim 4000$ eliminates this
        specific error.
  \item \textbf{\btt{maxHeartbeats} (200000).} Once the flags generate, elaborating
        the bridge theorem itself hits a \emph{deterministic timeout at} \btt{isDefEq}
        --- a different tunable. This is a \btt{native_decide}-on-a-huge-goal cost: the
        unifier works on the 816-element goal term independently of how fast the
        underlying computation is.
  \item \textbf{Wall-clock.} A reproduction with \btt{maxHeartbeats 0} (unlimited,
        chosen to ``measure true wall-clock'') and a \emph{non-killing} background
        build ran for $\sim 3$ hours without finishing --- because removing the
        heartbeat limit removed the only safety valve on the pathological
        \btt{isDefEq}. (Two process-management mistakes compounded here: a Bash
        \btt{timeout} parameter does \emph{not} terminate background tasks, and live
        \btt{lean.exe} processes seen during the run were the editor's LSP workers, not
        the build.)
\end{enumerate}
The lesson recorded in the history: \btt{maxRecDepth} is \emph{necessary but not
sufficient} for $n=6$; the genuine remaining wall is the bridge's elaboration-time
\btt{isDefEq} over an 800+-element goal, which none of the enumeration optimizations
touch.

\subsection{The \btt{unlabel} bug: a long red herring with a trivial root}

Separately, running \btt{generate_flags 3 0 6} produced, after $\sim$1--2 minutes, a
cryptic error:
\begin{lstlisting}
Application type mismatch: FlagAlgebras.flagEqv.refl ?m.14
  has type      ?m.14 ~f ?m.14
  but is expected to have type   ?m.9 = ?m.10
in   Quotient.sound (FlagAlgebras.flagEqv.refl ?m.14)
\end{lstlisting}
This is the proof of the generated \btt{unlabel_n_k_m_i} theorem
(\btt{unlabel Flag_n_k_m_i = Flag_n_0_0_j} by \btt{Quotient.sound (flagEqv.refl _)}).
A great deal of effort went into \emph{wrong} hypotheses: that it was an $n$-dependent
\emph{definitional-equality} failure (the underlying graphs not being defeq at $n=6$);
that it was a \btt{maxRecDepth} or \btt{maxHeartbeats} resource limit; that it needed
an iso-based proof instead of \btt{refl}. A minimal reproduction appeared to show
$n=5$ succeeding and $n=6$ failing with \emph{identical} code, ``confirming'' an
$n$-dependence --- but that minimal test was \emph{confounded}: it imported
\btt{FlagDef}, in which \btt{generate_empty_typed_flags 6} was commented out, so the
base flag \btt{Flag_6_0_0_0} did not exist and Lean silently \emph{auto-bound} it as an
opaque variable. Every ``defeq failure'' was really ``cannot prove
\btt{unlabel ... = (arbitrary variable)}''.

The actual root cause: \btt{generate_flags k m n} states each flag's \btt{unlabel}
bridge against the underlying empty-typed flag \btt{Flag_n_0_0_i}, which is produced by
\btt{generate_empty_typed_flags n}. With that prerequisite commented out, the bridge
references an undefined name, auto-bound as a variable, so the otherwise-correct proof
cannot close. Confirmation: with \btt{generate_empty_typed_flags 6} present, the exact
generated proof \btt{Quotient.sound (flagEqv.refl _)} compiles fine at $n=6$. So this
was never a math or performance bug at all --- it was a missing prerequisite surfaced
through a terrible error message.

\subsection{The fix}

Commit \btt{b15c62b} added a fast dependency check at the top of \btt{generate_flags}:
\begin{lstlisting}
unless (<- getEnv).contains (Name.mkSimple s!"Flag_{n}_0_0_0") do
  throwError s!"`generate_flags {k} {m} {n}` requires the underlying empty-typed flags
    `Flag_{n}_0_0_i`, ... Add `generate_empty_typed_flags {n}` before this command."
\end{lstlisting}
The cryptic, minutes-later \btt{flagEqv.refl} mismatch becomes an immediate,
actionable error (a 10-second build) when the prerequisite is absent, and the
with-prerequisite path is unchanged.

% =======================================================================
\section{Session work III: the degree-key dedup prefilter (\btt{54b06f7})}
\label{sec:degkey}
% =======================================================================

\subsection{The one untransferred technique}

After option~2, a systematic comparison of the two paths showed exactly one technique
present on the empty-typed side and absent on the labeled side: the
degree-sequence dedup prefilter (Section~\ref{sec:fbaeed0}). The labeled
\btt{dedupStepL} ran \btt{isIsoFast_bool} against \emph{every} survivor, whereas the
empty-typed \btt{dedupStepDeg} buckets by \btt{degKey}. Commit \btt{54b06f7} ports it.

The key choice was to key on the \emph{underlying graph's} \btt{degKey} rather than
invent a labeled-specific key:
\begin{lstlisting}
def labeledDegKey (G : Sym2LabeledGraph sigma n) := degKey <G.edges, G.edges_valid>
\end{lstlisting}
This is sound because a labeled $\sim_{\!sf}$ restricts to an $\sim_{\!sf}$ of the
underlying graphs (\btt{underlying_eqv_of_labeled_eqv}: take the same
\btt{graph_iso}, forget the empty \btt{type_preserve}), so
\btt{labeledDegKey_iso_invariant} follows directly from the empty-typed
\btt{degKey_iso_invariant}. The fold simulation \btt{foldl_dedupStepDegL_sim} is a
near-verbatim copy of \btt{foldl_dedupStepDeg_sim}, and \btt{genLabeledGraphsDedup_eq}
(the analogue of \btt{augRepsDeg_fst_eq}) proves the keyed dedup \emph{equals} the
naive \btt{dedupStepL} fold --- so the value, the bridge, and all
completeness/\btt{Nodup} proofs are unchanged (they were re-routed through the
equality). Keying on the underlying graph captures the bulk of the benefit: a flag of
underlying graph $G$ then runs the expensive iso test only against survivors whose
underlying graph shares $G$'s degree key --- almost always just survivors with the
\emph{same} underlying graph, rather than all $\sim g$ of them.

\subsection{An anticipated improvement that did not materialize}

This was flagged in advance as ``a real untransferred technique, but probably not the
$n=6$ bottleneck'', and the measurement confirmed it: the dedup over the post-option~2
2,808-element list is already \emph{sub-10\,ms} either way. The reasons are structural
and were predicted:
\begin{itemize}
  \item Option~2 had already shrunk the dedup input $175\times$, so the fold is small.
  \item The labeled \btt{isIsoFast_bool} is \emph{already} cheap: it prefilters on
        edge count, and --- because a flag isomorphism must preserve the type embedding
        --- it permutes only the $(n-k)!$ \emph{non-type} vertices, i.e.\ $3!=6$ at
        $n=6,\,k=3$, not $n!=720$. The empty-typed \btt{degKey} mattered precisely
        because the unlabeled check did pay $n!$; on the labeled side that motivation
        is largely gone.
\end{itemize}
The change is nonetheless correct, \btt{sorry}-free, and value-preserving, and was kept
(by explicit user decision) for structural parity with the empty-typed path and as
insurance if the dedup ever becomes a bottleneck at larger $n$. The honest accounting:
it optimizes a step that option~2 had already made cheap, and it does \emph{not}
address the actual $n=6$ wall (the bridge \btt{isDefEq}). The
\btt{naive\_ms=0\ keyed\_ms=0} reading from the in-\btt{\#eval} timer was itself partly
a measurement artifact (Lean \btt{let} bindings are lazy, so the work happened after the
timing points), but the build-profile-level conclusion --- ``already fast either way''
--- is robust.

% =======================================================================
\section{Complexity of the successive approaches}
% =======================================================================

The table summarizes how the two halves of each command scaled across the history.
Here $g$ is the number of flags on a line, $R$ the size of the \emph{raw} candidate
list fed to the dedup, and ``iso check'' the cost of one $\sim_{\!sf}$ decision.

\begin{center}
\begin{tabular}{lll}
\toprule
component & before & after \\
\midrule
empty-typed enumeration   & augmentation (always) & augmentation $+$ degree prefilter \\
labeled raw list $R$ ($n{=}6$) & $491{,}520$ (brute force) & $2{,}808$ (augmentation) \\
labeled dedup             & $O(R\cdot g)$ iso checks, all-pairs & bucketed by \btt{degKey} \\
empty-typed completeness  & $O(g^2)$ \btt{toFinset} bridge & $O(g)$ positional bridge \\
labeled completeness      & $O(g^2)$ \btt{Finset} $+$ $O(g^2)$ \btt{Nodup} & $O(g)$ positional $+$ proof \\
\btt{= univ}              & monster \btt{Fintype} \btt{native\_decide} & symbolic theorem $+$ cheap bridge \\
per-flag \btt{type\_embed}  & \btt{fin\_cases <;> decide} (per-case instance) & whole-goal \btt{decide} \\
per-flag \btt{edges\_valid} & \btt{by decide} (kernel re-reduction) & structural \btt{fin\_cases} \\
\bottomrule
\end{tabular}
\end{center}

The most important architectural point is that two \emph{separate} ideas were needed
to make the labeled completeness tractable, and they compose: the
\emph{augmentation enumeration} (commit \btt{a993a33}) makes the thing the bridge
\emph{evaluates} cheap, while the \emph{$O(g)$ ordered bridge} (commit \btt{bb6f60f})
makes the \emph{comparison} cheap. Neither alone suffices; the brute-force enumeration
under an $O(g)$ bridge is still exponential, and an augmentation enumeration under an
$O(g^2)$ bridge still pays the quadratic comparison.

% =======================================================================
\section{Measured and estimated speedups (summary)}
% =======================================================================

\begin{center}
\begin{longtable}{p{2.1cm}p{5.4cm}p{6.0cm}}
\toprule
commit & change & measured / estimated effect \\
\midrule
\endhead
\btt{212968a} & Phase 2: typed Lean generation $+$ symbolic completeness & \btt{FlagDef} build $568\,\text{s}\to274\,\text{s}$ ($\sim 2\times$) \\
\btt{fbaeed0} & degree prefilter, decorate-sort, $O(g)$ empty-typed bridge, $n{=}6$ & enabled $n{=}6$; $O(g^2)\to O(g)$ empty-typed bridge \\
\btt{fa8939e} & structural \btt{edges_valid} & $n{=}6$ \btt{decide} $\sim 9\times$; kernel type-checking $\sim$ halved \\
\btt{de7dcbc} & reuse precomputed \btt{canonicalEdgeList} & $n{=}6$ empty-typed $13.6\,\text{s}\to9.73\,\text{s}$ ($\sim 28\%$) \\
\btt{81cf849} & whole-goal \btt{decide} for \btt{type_embed} & heaviest line $\sim 63\,\text{s}\to22\,\text{s}$; \btt{FlagDef} $\sim 301\,\text{s}\to100\,\text{s}$ ($\sim 3\times$) \\
\btt{bb6f60f} & $O(g)$ labeled bridge (this session) & per-line $17.6\,\text{s}\to7.7\,\text{s}$ at $n{=}5$ ($\sim 2.3\times$); removes the $13\,\text{s}$ \btt{Nodup} \btt{native_decide} \\
\btt{a993a33} & augmentation labeled enumeration (this session) & raw list $491{,}520\to2{,}808$ at $n{=}6$ ($175\times$); dedup $\sim 10^8\to 2.3\times10^6$ iso checks \\
\btt{b15c62b} & prerequisite check (this session) & cryptic minutes-later error $\to$ clear $10$-second error \\
\btt{54b06f7} & labeled degree-key prefilter (this session) & negligible at $n{=}6$ (dedup already sub-$10$\,ms); structural parity \\
\bottomrule
\end{longtable}
\end{center}

% =======================================================================
\section{Remaining bottlenecks and what became impractical}
% =======================================================================

\subsection{$n=6$ (labeled): the bridge \btt{isDefEq} wall}

After every optimization above, \btt{generate_flags 3 0 6} is still not a routine
build. The remaining costs, in rough order:
\begin{itemize}
  \item The bridge \btt{native_decide}'s \emph{elaboration-time} \btt{isDefEq} over the
        816-element goal term --- this is what times out at the default
        \btt{maxHeartbeats}, and it is \emph{not} addressed by any enumeration or
        comparison optimization, because it is about the syntactic size of the goal,
        not the cost of the computation.
  \item The 816 per-flag \btt{type_embed} \btt{decide}s (already $\sim 3\times$ cheaper
        after \btt{81cf849}, but $11\times$ more of them than at $n=5$).
  \item The depth-816 list literals, which require \btt{set_option maxRecDepth 4000}.
\end{itemize}
None of these is the dedup, which option~2 and the degree-key prefilter addressed.
This is why the degree-key prefilter, although the ``textbook'' missing optimization,
did not move the needle: it sharpened a part that was no longer the bottleneck.

\subsection{$n=7$ (empty-typed): the interpreter-enumeration wall}

$n=7$ (1044 classes) is omitted entirely. Two independent reasons are documented in
\btt{FlagDef.lean}: the elaboration-time enumeration runs in the Lean interpreter and
takes $>80$ minutes, and the 1044-element list literal in the completeness bridge
exceeds the default \btt{maxRecDepth}. Both are properties of \emph{interpreter-based
generation}, not of the mathematics; the file note explicitly says ``revisit if the
generator is moved to compiled (native) evaluation''.

\subsection{What is genuinely hard to optimize further, and why}

\begin{itemize}
  \item The \btt{native_decide} \emph{kernel evaluation} (\btt{reduceBool}) is, by
        design, the trust boundary; its cost is the cost of running the decision
        procedure once and cannot be ``proved away''. The only levers are making the
        procedure cheaper (done: augmentation, $(n-k)!$ search, prefilters) or making
        the \emph{goal} smaller (not yet done).
  \item The large-goal \btt{isDefEq} is a property of stating a single equality between
        an $\sim g$-element literal and a computed list. Shrinking it means changing
        the \emph{shape} of the bridge (Section~\ref{sec:future}), not tuning a tactic.
  \item The empty-typed path is essentially at the limit of the current strategy: it
        already uses canonical augmentation, a degree prefilter, decorate-sort, key
        reuse, structural \btt{edges_valid}, and an $O(g)$ bridge. The only remaining
        large win there is the move off the interpreter.
\end{itemize}

% =======================================================================
\section{Dead ends, false starts, and reverted experiments}
% =======================================================================

\begin{itemize}
  \item \textbf{Recomputing the sort key in the comparator.} The first
        \btt{genFlagsOrdered} recomputed the $n!$-scanning \btt{labeledFlagKey} inside
        the sort comparator, making it $O(g^2)$ key computations; replaced by
        decorate-sort-undecorate.
  \item \textbf{The ``new bridge is slower'' near-miss.} Comparing \btt{measO1}
        ($7.7$\,s) to \btt{measO2} ($4.5$\,s) alone suggested a regression; the fair
        comparison (against \btt{measO2}$+$\btt{measNodup}$=17.6$\,s) reversed the
        conclusion.
  \item \textbf{\btt{maxRecDepth} / \btt{maxHeartbeats} as the $n=6$ fix.} Both were
        tried and both only \emph{revealed the next wall}; neither makes $n=6$ a
        routine build.
  \item \textbf{The 3-hour runaway.} \btt{maxHeartbeats 0} plus a misunderstanding that
        a background build would be auto-killed produced a multi-hour, ultimately
        useless run.
  \item \textbf{The \btt{unlabel} ``defeq bug''.} A long investigation chased an
        $n$-dependent definitional-equality failure that did not exist; the minimal
        reproduction was confounded by an auto-bound \btt{Flag_6_0_0_0} (because the
        imported \btt{FlagDef} had \btt{generate_empty_typed_flags 6} commented out).
        The real cause was a missing prerequisite (Section~\ref{sec:n6frontier}).
  \item \textbf{The degree-key prefilter as a speedup.} Implemented and proved, but the
        measured benefit at $n=6$ was negligible; kept for structural parity rather
        than performance (Section~\ref{sec:degkey}).
\end{itemize}

% =======================================================================
\section{Promising directions for future optimization}
\label{sec:future}
% =======================================================================

\begin{enumerate}
  \item \textbf{Compiled (native) elaboration-time generation.} The single highest-
        value change. \btt{genFlagData}/\btt{genCanonicalEdgeLists} are already
        evaluated via compiled \btt{Lean.Meta.evalExpr} for the \emph{data}, but the
        broader enumeration and the bridges still hit interpreter and kernel-reduction
        costs. Moving the heavy enumeration fully onto compiled evaluation is what the
        $n=7$ note explicitly anticipates, and it directly attacks the $>80$-minute
        interpreter wall.
  \item \textbf{Shrinking the bridge goal.} The $n=6$ wall is the \btt{isDefEq}/
        \btt{native_decide} over an $\sim g$-element literal. Emitting the list as a
        \emph{chunked} concatenation (each chunk under \btt{maxRecDepth}, with a shallow
        \btt{++} spine) would avoid the depth-$g$ literal and shrink the per-step goal,
        at the cost of a more elaborate bridge proof. This is the only idea that
        targets the actual remaining wall, but it touches every bridge proof and was
        deliberately deferred.
  \item \textbf{Structural per-flag \btt{type_embed} proofs.} The 816 per-flag
        \btt{decide}s are the second-largest $n=6$ cost and are \emph{typed-specific}
        (the empty-typed path has no type embedding, so there is nothing to port). The
        \btt{edges_valid} refactor (Section~\ref{sec:edges_valid}) is the template: a
        compact structural proof in place of a kernel \btt{decide}.
  \item \textbf{A type-refined degree key.} The labeled \btt{degKey} currently ignores
        the type placement. Adding the ordered sequence of type-vertex degrees
        $[\deg(\btt{type_embed}\ i)]_i$ --- itself iso-invariant via
        \btt{type_preserve} --- would split within-underlying-graph buckets further.
        The expected gain is small (within-graph buckets are already tiny), so this is
        a refinement, not a fix.
  \item \textbf{Caching canonical forms per underlying graph in \btt{labeledFlagKey}.}
        The sort key recomputes the $n!$ canonical form per \emph{flag}, whereas
        \btt{genFlagData} computes it once per \emph{underlying graph}. Reusing it
        (as \btt{genCanonicalEdgeLists} did for the empty-typed sort,
        Section~\ref{sec:reuse}) would cut the key-computation part of the bridge's
        native run --- though, again, the native run is not the $n=6$ wall.
\end{enumerate}

% =======================================================================
\section{Conclusion}
% =======================================================================

Over this history the flag generators moved from an external JSON pipeline to a fully
self-contained Lean development, and the dominant proof obligation --- completeness ---
went from a prohibitive ``monster'' \btt{Fintype} \btt{native_decide} to a symbolic
theorem bridged by a single $O(g)$ positional comparison. The two structurally
significant wins were internalizing the typed enumeration with a symbolic completeness
proof (\btt{212968a}, $\sim 2\times$) and, this session, giving the labeled path the two
techniques the empty-typed path already had: an augmentation enumeration in place of a
$2^{\binom n2}$ brute force (\btt{a993a33}, $175\times$ fewer raw graphs at $n=6$) and a
linear ordered completeness bridge in place of two quadratic \btt{native_decide}s
(\btt{bb6f60f}, $\sim 2.3\times$ per line and the removal of the $13$-second \btt{Nodup}
check). A series of proof-engineering optimizations --- structural \btt{edges_valid}
($\sim 9\times$ \btt{decide}), whole-goal \btt{decide} ($\sim 3\times$ on the full
build), and canonical-form reuse ($\sim 28\%$) --- compounded with these.

The honest closing assessment is that the development is now bounded not by the
enumeration or the comparison, both of which are cheap, but by two costs that the
\emph{shape} of the current architecture imposes: the interpreter-time enumeration
(fatal at $n=7$) and the elaboration-time \btt{isDefEq}/\btt{native_decide} over a
$g$-element goal literal (expensive at $n=6$). The one optimization that the empty-typed
path had and the labeled path lacked --- the degree-key dedup prefilter --- was duly
ported, proved correct, and found to be negligible in practice, precisely because the
earlier augmentation change had already removed the cost it was meant to address. The
clearest path forward is compiled generation and a chunked bridge goal; everything else
is refinement.

\end{document}
